\documentclass[11pt]{article}
\usepackage[margin=1in]{geometry}
\usepackage{amsmath,amssymb,amsthm}
\usepackage{graphicx}
\usepackage[hidelinks]{hyperref}

\newtheorem{theorem}{Theorem}
\newtheorem{lemma}{Lemma}
\newtheorem{corollary}{Corollary}

\newcommand{\ellone}{\ell_1}
\newcommand{\ellinf}{\ell_\infty}
\newcommand{\czero}{c_0}
\newcommand{\wuc}{weakly unconditionally Cauchy}

\begin{document}

\begin{center}
{\Large Property-\((u)\) \(\ell_1\)-preduals must contain \(c_0\)}\\[0.5em]
{\small Run \texttt{fa\_banach\_001}, arXiv:1504.08223, candidate partial result}
\end{center}

\section*{Source and Question}

The source paper is S. A. Argyros, I. Gasparis, and P. Motakis,
\emph{On the structure of separable \(\mathcal L_\infty\)-spaces},
arXiv:1504.08223 \cite{AGM}.  In the introduction, the authors describe
their asymptotic \(\czero\) isomorphic \(\ellone\)-predual without \(\czero\)
as a step toward Godefroy--Li Question IV.2:

\begin{center}
\includegraphics[width=\linewidth]{figures/open_problem_crop.png}
\end{center}

The relevant source question is whether every isomorphic \(\ellone\)-predual
satisfying Pe\l czy\'nski's property-\((u)\) is isomorphic to \(\czero\).

\section*{Result}

\textbf{Status.}  This packet proves a partial obstruction, not the full
Godefroy--Li classification.  It shows that property-\((u)\) forces the
presence of \(\czero\) in every infinite-dimensional isomorphic \(\ellone\)-
predual.  The remaining question is whether the whole space must be
isomorphic to \(\czero\).

\section*{Definitions}

We use the standard formulation of Pe\l czy\'nski's property-\((u)\): for every
weakly Cauchy sequence \((x_n)\) in \(X\), there is a \wuc{} series
\(\sum_{k=1}^\infty u_k\) in \(X\) such that
\[
  x_n-\sum_{k=1}^n u_k \longrightarrow 0
\]
weakly.  Recall also the Bessaga--Pe\l czy\'nski theorem: a Banach space
contains no isomorphic copy of \(\czero\) if and only if every \wuc{} series in
it is unconditionally norm convergent \cite{AK, BP}.

\section*{Proof Intuition}

Property-\((u)\) says that weakly Cauchy sequences are controlled by partial
sums of \wuc{} series.  The only obstruction to turning such partial sums into
actual norm-convergent sums is \(\czero\): without \(\czero\), the
Bessaga--Pe\l czy\'nski theorem makes every \wuc{} series unconditionally
convergent.  Thus property-\((u)\) plus no \(\czero\) upgrades weak Cauchy
sequences to weakly convergent sequences.

For an \(\ellone\)-predual, the dual is separable, so the space cannot contain
\(\ellone\).  Rosenthal's \(\ellone\)-theorem then says that any nonreflexive
bounded sequence has a weakly Cauchy subsequence; weak sequential completeness
would make that subsequence weakly convergent, contradicting
Eberlein--Smulian.  Hence no-\(\czero\) would force reflexivity, impossible
for an infinite-dimensional space whose dual is \(\ellone\).

\begin{lemma}\label{lem:u-no-c0-wsc}
If \(X\) has property-\((u)\) and contains no isomorphic copy of \(\czero\),
then \(X\) is weakly sequentially complete.
\end{lemma}

\begin{proof}
Let \((x_n)\) be a weakly Cauchy sequence in \(X\).  By property-\((u)\), there
is a \wuc{} series \(\sum u_k\) such that
\[
  x_n-s_n \to 0 \quad\text{weakly}, \qquad
  s_n=\sum_{k=1}^n u_k .
\]
Since \(X\) contains no \(\czero\), the Bessaga--Pe\l czy\'nski theorem implies
that \(\sum u_k\) is unconditionally norm convergent.  Hence \(s_n\to s\) in
norm for some \(s\in X\).  Therefore \(x_n=(x_n-s_n)+s_n\to s\) weakly.
\end{proof}

\begin{theorem}\label{thm:main}
Let \(X\) be an infinite-dimensional Banach space such that \(X^*\) is
isomorphic to \(\ellone\).  If \(X\) has Pe\l czy\'nski's property-\((u)\),
then \(X\) contains an isomorphic copy of \(\czero\).
\end{theorem}

\begin{proof}
Assume toward a contradiction that \(X\) contains no copy of \(\czero\).
By Lemma~\ref{lem:u-no-c0-wsc}, \(X\) is weakly sequentially complete.

Since \(X^*\simeq \ellone\), the dual \(X^*\) is separable.  Thus \(X\)
contains no copy of \(\ellone\): if \(J:\ellone\to X\) were an isomorphic
embedding, then \(J^*:X^*\to \ellinf\) would be a bounded surjection, making
\(\ellinf\) a separable quotient of the separable space \(X^*\), a
contradiction.

We claim that \(X\) is reflexive.  If not, the closed unit ball of \(X\) is not
weakly compact.  By Eberlein--Smulian there is a bounded sequence in \(X\) with
no weakly convergent subsequence.  Since \(X\) contains no \(\ellone\),
Rosenthal's \(\ellone\)-theorem gives a weakly Cauchy subsequence.  Weak
sequential completeness makes this subsequence weakly convergent, a
contradiction.  Hence \(X\) is reflexive.

But if \(X\) is reflexive, then \(X^*\) is reflexive.  This is impossible
because \(X^*\simeq \ellone\) and infinite-dimensional \(\ellone\) is not
reflexive.  Therefore the assumption that \(X\) contains no \(\czero\) was
false.
\end{proof}

\begin{corollary}[Complemented \(\czero\) self-reduction]\label{cor:self}
Let \(X\) satisfy the hypotheses of Theorem~\ref{thm:main}.  Then
\[
  X \simeq \czero \oplus Y
\]
for a complemented subspace \(Y\) of \(X\).  Moreover, either \(Y\) is
finite-dimensional, in which case \(X\simeq \czero\), or \(Y\) is again an
infinite-dimensional Banach space with property-\((u)\) and
\(Y^*\simeq \ellone\).
\end{corollary}

\begin{proof}
By Theorem~\ref{thm:main}, \(X\) contains a subspace \(G\simeq \czero\).
Since \(X^*\simeq \ellone\), the space \(X\) is separable.  Sobczyk's theorem
therefore gives a bounded projection \(P:X\to G\).  Put \(Y=\ker P\).  Then
\(X\simeq G\oplus Y\simeq \czero\oplus Y\).

If \(Y\) is finite-dimensional, then \(\czero\oplus Y\simeq \czero\).  Assume
that \(Y\) is infinite-dimensional.  Property-\((u)\) passes to complemented
subspaces: if \((y_n)\) is weakly Cauchy in \(Y\), apply property-\((u)\) in
\(X\) and then apply the projection \(I-P\) to the witnessing \wuc{} series.
Thus \(Y\) has property-\((u)\).

Finally, \(Y^*\) is isomorphic to a complemented subspace of
\(X^*\simeq \ellone\).  By Pe\l czy\'nski's theorem on complemented subspaces
of \(\ellone\), every infinite-dimensional complemented subspace of
\(\ellone\) is isomorphic to \(\ellone\) \cite{LT}.  Hence
\(Y^*\simeq \ellone\).
\end{proof}

\section*{Verification and Novelty Check}

The proof is purely structural and uses no computation.  The cheap run indexes
were searched for arXiv:1504.08223 and for the phrases
``property-(u)'', ``Question IV.2'', and isomorphic \(\ellone\)-predual; no
existing packet for this exact partial theorem was found.  A nearby packet on
the Rosenthal--Odell quotient problem records a broader isomorphic
\(\ellone\)-predual counterexample context, but not the property-\((u)\)
obstruction proved here.  A bounded web/local literature check on
2026-07-03 located standard definitions and the Bessaga--Pe\l czy\'nski/Rosenthal
tools, but no later exact answer to the full Godefroy--Li isomorphism question.

\section*{Limitations}

The theorem and Corollary~\ref{cor:self} do not prove that \(X\) is isomorphic
to \(\czero\).  They rule out the possibility that a counterexample to
Godefroy--Li Question IV.2 could avoid \(\czero\) entirely, and they show that
any remaining counterexample recursively splits off complemented copies of
\(\czero\).  The missing step is uniform control or a structural argument that
forces the complemented remainders to terminate or be absorbed by \(\czero\).

\section*{Human Review Recommendation}

Review as a likely valid partial result.  The main checks are the exact
definition of property-\((u)\), the use of the Bessaga--Pe\l czy\'nski
characterization of spaces with no \(\czero\), and the separable-dual
exclusion of \(\ellone\) subspaces.

\begin{thebibliography}{9}

\bibitem{AGM}
S. A. Argyros, I. Gasparis, and P. Motakis,
\emph{On the structure of separable \(\mathcal L_\infty\)-spaces},
arXiv:1504.08223, 2015.

\bibitem{AK}
F. Albiac and N. J. Kalton,
\emph{Topics in Banach Space Theory},
Graduate Texts in Mathematics 233, Springer, 2006.

\bibitem{BP}
C. Bessaga and A. Pe\l czy\'nski,
\emph{On bases and unconditional convergence of series in Banach spaces},
Studia Math. \textbf{17} (1958), 151--164.

\bibitem{Rosenthal}
H. P. Rosenthal,
\emph{A characterization of Banach spaces containing \(\ell_1\)},
Proc. Nat. Acad. Sci. U.S.A. \textbf{71} (1974), 2411--2413.

\bibitem{LT}
J. Lindenstrauss and L. Tzafriri,
\emph{Classical Banach Spaces I: Sequence Spaces},
Ergebnisse der Mathematik und ihrer Grenzgebiete 92, Springer, 1977.

\end{thebibliography}

\end{document}
